\section{Introducción}

\subsection{Contexto y motivación}

Las asociaciones civiles deportivas sin fines de lucro en la República Argentina constituyen pilares fundamentales para la cohesión social, la contención comunitaria y la promoción de actividades culturales y recreativas. A diferencia del modelo societario de propiedad privada imperante en otras latitudes, el club deportivo argentino opera tradicionalmente bajo un esquema democrático donde el socio es el verdadero propietario y gestor de la institución. No obstante, a pesar de su incalculable valor sociocultural, la inmensa mayoría de estas entidades de tamaño mediano y pequeño ---desde los clubes de las categorías de ascenso del área metropolitana de Buenos Aires hasta las ligas regionales del interior del país--- enfrentan un escenario de marcada vulnerabilidad financiera, profundizado por la subsistencia de prácticas de gestión administrativa marcadamente tradicionales y reactivas.

Históricamente, la adopción de herramientas tecnológicas avanzadas ha sido un privilegio exclusivo de las escasas corporaciones deportivas que de forma sistemática acceden a los recursos de la élite profesional de primera división. Para las instituciones del ascenso y del interior, la digitalización de sus operaciones suele limitarse a la implementación de planillas de cálculo manuales aisladas, sistemas de base de datos desactualizados o simples sistemas de registro de transacciones genéricos de tipo reactivo, conocidos técnicamente como sistemas \textit{CRUD} básicos (crear, leer, actualizar y borrar). Esta brecha tecnológica genera un aislamiento operativo que restringe la capacidad de las dirigencias para automatizar procesos administrativos rutinarios, optimizar la recaudación de cuotas, prevenir fugas de capital y garantizar la seguridad perimetral de sus instalaciones durante eventos de concurrencia masiva.

Frente a esta realidad, el presente Trabajo Profesional documenta el diseño, desarrollo e implementación de \textbf{SocioUnido}, un ecosistema de \textit{software} modular concebido bajo un modelo de \textit{software} como servicio (\textit{SaaS}) y marca blanca, específicamente dimensionado para resolver las demandas operativas reales de las instituciones deportivas argentinas. El proyecto no se concibe como una simple herramienta contable digital, sino como una plataforma inteligente, descentralizada y proactiva de inteligencia de negocios y ciberseguridad aplicada al saneamiento financiero y la inclusión tecnológica institucional.

\subsection{Planteamiento del problema y faltante en el sector}

El análisis de dominio realizado sobre la gestión operativa de los clubes deportivos en Argentina permitió identificar tres ineficiencias críticas de carácter silencioso que impactan directamente sobre su rentabilidad económica y su sostenibilidad institucional:

\begin{enumerate}
    \item \textbf{La morosidad societaria silenciosa}: Debido a la desconexión existente entre las actividades deportivas de los socios, los registros administrativos y los medios de pago disponibles, las instituciones deportivas sufren una pérdida constante de ingresos por cuotas sociales impagas. La falta de canales digitales inmediatos y la dependencia de cobros manuales presenciales en tesorerías centralizadas fomentan un comportamiento de olvido por parte de los usuarios. Al no contar con alertas proactivas ni motores capaces de anticipar la morosidad o la intención de baja societaria, los clubes actúan reactivamente una vez que la deuda se ha consolidado y la pérdida del socio es irreversible.
    \item \textbf{El fraude en el control de accesos por duplicación o préstamo de credenciales físicas}: En eventos masivos, la utilización de credenciales impresas plásticas con códigos de barras estáticos o sistemas magnéticos heredados propicia el préstamo de carnets entre socios activos y personas no autorizadas. Esta duplicidad de accesos genera no solo una severa fuga de recaudación por entradas no vendidas, sino también una grave vulnerabilidad de seguridad lógica y física en los puntos de control del estadio. La solución tradicional exige conectividad permanente de banda ancha de alta velocidad en los molinetes perimetrales para validar la validez de los códigos contra la base de datos central en la nube. Sin embargo, en escenarios masivos, la saturación por radiofrecuencia y la caída recurrente de las redes de telefonía celular móvil imposibilitan la validación en tiempo real (\textit{online}), obligando en muchas ocasiones al personal de control a liberar los ingresos para evitar desmanes físicos, validando involuntariamente accesos fraudulentos o de socios morosos.
    \item \textbf{La subutilización de instalaciones comunes, disciplinas y activos deportivos rentables}: Gran parte de los clubes poseen instalaciones ociosas (como canchas auxiliares, quinchos, salones o predios de entrenamiento) e infraestructura de disciplinas deportivas infrautilizadas debido a la ausencia de un canal centralizado de reservas. Los socios deben acudir presencialmente o comunicarse telefónicamente en horarios acotados para consultar disponibilidad, lo que desincentiva el consumo interno y reduce las fuentes de financiamiento alternativas de la institución.
\end{enumerate}

SocioUnido nace con el firme propósito de mitigar estas problemáticas de raíz, integrando en una única plataforma la gestión administrativa y social proactiva, la prevención analítica de la morosidad y un sistema de control de accesos seguro y descentralizado capaz de validar ingresos masivos de forma 100\% fuera de línea (\textit{offline}).

\subsection{Objetivos del Trabajo Profesional}

\subsubsection{Objetivo general}
Diseñar, desarrollar y validar un ecosistema de \textit{software} integral y adaptativo de marca blanca que centralice la administración de clubes deportivos, optimice la retención y fidelización societaria mediante mecanismos cognitivos y predictivos, y blinde el perímetro de las instalaciones mediante una solución criptográfica de control de accesos descentralizada y tolerante a fallos de conectividad.

\subsubsection{Objetivos específicos}
Para alcanzar el objetivo general, se definen las siguientes metas técnicas y metodológicas:
\begin{itemize}
    \item Implementar una arquitectura perimetral y distribuida de microservicios resiliente, escalable y cibersegura, orquestada mediante contenedores y centralizada bajo un portal de enlace de aplicaciones (\textit{API gateway}) unificado.
    \item Desarrollar una aplicación web progresiva (\textit{PWA}) adaptativa para el socio, compatible con dispositivos móviles de bajos recursos, que centralice su credencial digital, la reserva de espacios, la inscripción a disciplinas y el flujo transaccional de pagos.
    \item Diseñar y validar un algoritmo de control de accesos inteligente basado en códigos criptográficos de un solo uso basados en tiempo (\textit{TOTP}), que permita la validación instantánea del estado de pago e identidad del socio sin requerir conectividad de red (\textit{offline}) en los dispositivos de escaneo del club.
    \item Construir un asistente transaccional conversacional interactivo cognitivo (asistente ``BotIn'') montado sobre la API de Telegram, con procesamiento de lenguaje natural (\textit{PLN}), para reducir la fricción transaccional administrativa a cero.
    \item Diseñar y desarrollar una plataforma web de administración centralizada que permita a las autoridades del club digitalizar de forma integral los procesos institucionales, automatizar la facturación y cobranza de cuotas sociales, y gestionar de manera eficiente las disciplinas deportivas y los activos compartidos.
    \item Validar metodológicamente el proceso de ingeniería de \textit{software} utilizando el marco ágil Scrum y certificar la calidad, ciberseguridad e interoperabilidad del código fuente de forma continua mediante el análisis estático y un agente autónomo de aseguramiento de calidad (\textit{QA}) basado en inteligencia artificial.
\end{itemize}

\subsection{Articulación e integración de conocimientos académicos de la carrera}

El diseño y materialización del ecosistema SocioUnido no constituye una actividad de desarrollo tecnológico aislada, sino la síntesis práctica y la convergencia multidisciplinar de los conocimientos teóricos adquiridos a lo largo de toda la trayectoria académica en la Facultad de Ingeniería de la Universidad de Buenos Aires. La resolución de los desafíos del dominio deportivo y de infraestructura requirió articular de forma sinérgica conocimientos de áreas diversas de la informática y de la formación general del ingeniero:

\begin{itemize}
    \item \textbf{Ciencias básicas e infraestructura física}: Los fundamentos de electromagnetismo, atenuación de ondas de radiofrecuencia y saturación de espectro analizados en las asignaturas de Física y Física para Informática proveyeron el sustento científico para predecir el colapso de las antenas celulares en los estadios de fútbol ante altas densidades de usuarios. Esta base física justificó de manera irrefutable la necesidad de diseñar un sistema de control de ingresos descentralizado y desacoplado de la red en la nube.
    \item \textbf{Matemáticas avanzadas y criptografía}: Las teorías de aritmética modular, funciones unidireccionales y transformaciones matriciales estructuradas en Álgebra y Álgebra Lineal permitieron desarrollar y robustecer el algoritmo criptográfico del sistema de accesos \textit{offline} mediante semillas \textit{TOTP} e implementar firmas digitales con claves simétricas seguras e infalsificables.
    \item \textbf{Ingeniería de software, diseño y arquitectura distribuidos}: Los conceptos de encapsulamiento, modularidad y diseño orientado a objetos (entre muchos otros) provistos por Paradigmas de Programación y Taller de Programación guiaron la construcción de un código fuente limpio y mantenible. Por su parte, Base de Datos aportó el marco teórico del modelo relacional, control de transacciones concurrentes complejas y optimización de consultas SQL en PostgreSQL. Redes y Sistemas Distribuidos I fueron la base técnica para diseñar la topología de microservicios, la tolerancia a fallos y la comunicación interservicio asíncrona.
    \item \textbf{Ciberseguridad y defensa profunda}: La aplicación de los lineamientos de análisis de vulnerabilidades, sanitización de entradas, prevención de inyecciones y control estricto de roles de acceso mediante JSON Web Tokens (\textit{JWT}) provienen directamente del Taller de Seguridad Informática, construyendo una infraestructura segura contra ataques perimetrales.
    \item \textbf{Emprendedurismo, ética y contexto profesional}: El desarrollo de un producto mínimo viable comercializable, el estudio macrocontextual bajo la metodología de las 5 C, el análisis microcontextual de las 4 P, el seguimiento de métricas \textit{SaaS} y la gestión ágil \textit{Scrum} se forjaron de manera directa sobre las bases conceptuales y procedimentales adquiridas a lo largo de un trayecto de cuatro materias clave: ingeniería del software I, ingeniería del software II, empresas de base tecnológica I y empresas de base tecnológica II. Estas asignaturas proveyeron el andamiaje metodológico para estructurar tanto el ciclo de vida del producto de software como su estrategia de viabilidad y comercialización en el mercado real.
\end{itemize}

Cabe destacar que en el capítulo de conclusiones se realiza un desglose más exhaustivo y pormenorizado acerca de cómo cada una de estas asignaturas influyó de manera directa en el proceso de diseño, desarrollo y validación del presente trabajo profesional.
